% Methods — Paper 2, Medical Physics submission
% Humanised v1, 2026-05-25, formal-academic-active register.
% Every methodological claim is bound to a citation in references.bib.
% Draft → self-critique → revision pass per /humanizer skill protocol.

\subsection{Study design and pre-registration}

This work is a prospectively designed, controlled-perturbation in-silico reproducibility study comparing the propagation of segmentation-mask perturbation through pretrained foundation-model (FM) embeddings and IBSI-aligned handcrafted radiomic features. All imaging data are secondary, de-identified, and publicly released; the analyses use the Lung Image Database Consortium and Image Database Resource Initiative (LIDC-IDRI) collection~\cite{armato2011} hosted by The Cancer Imaging Archive (TCIA)~\cite{clark2013}. The study qualifies for ethics exemption under the standard TCIA Data Use Agreement and was conducted without new patient contact. The analytical protocol, including four pre-registered hypotheses, inclusion criteria, perturbation specification, statistical analysis plan, and target journal, was deposited on the Open Science Framework (OSF DOI: [pending]) prior to any feature extraction. Protocol amendments necessitated during execution were recorded in a versioned decisions log; the substantive amendments are summarised in this Methods section, and the full log is provided as Supplementary File 1.

\subsection{Dataset and cohort}

The LIDC-IDRI collection comprises 1\,018 thoracic computed-tomography scans with annotations from up to four independent expert radiologists per nodule, including binary segmentation masks and a 1--5 ordinal malignancy rating~\cite{armato2011}. The inclusion criteria specified nodules with (i) segmentations from at least three readers, (ii) median diameter at or above 6~mm --- the Fleischner Society 2017 actionable threshold for solid pulmonary nodules~\cite{macmahon2017} --- and (iii) median malignancy rating $\leq 2$ (benign) or $\geq 4$ (malignant), with intermediate ratings excluded. The 6~mm floor was reached via a pre-registered fallback cascade after the original 10~mm floor yielded only 29 benign nodules across 600 downloaded patients; the corresponding sensitivity analysis at 10~mm is reported in Supplementary Table~S2.

Two hundred LIDC-IDRI series were initially downloaded from TCIA, with a further 400 added to obtain sufficient benign-class nodules, for a total of 600 series ($\sim$52~GB). Annotation parsing was performed with pylidc 0.2.3~\cite{hancock2016pylidc}, and reader segmentations were aligned within a unified bounding box per nodule using the \texttt{pylidc.utils.consensus} routine. The inclusion cascade produced 353 nodules (198 malignant, 155 benign) from 251 patients. A balanced cohort was constructed by seeded random downsampling of the majority class to the minority count, producing the primary analysis cohort of 310 nodules (155 malignant, 155 benign) from 227 unique patients. The unbalanced 353-nodule cohort was retained for sensitivity analyses.

Cohort descriptive statistics: benign nodules had median diameter 7.74~mm (interquartile range, IQR, 6.70--9.20~mm; range 6.00--35.36~mm); malignant nodules had median diameter 19.36~mm (IQR 14.30--25.58~mm; range 6.47--44.60~mm). All nodules satisfied the $\geq 3$-reader threshold and 65\% had segmentations from all four readers. The median pairwise inter-reader Dice similarity coefficient across the cohort was 0.789 (mean 0.769, range 0.156--0.934), consistent with previously reported LIDC inter-reader Dice values of 0.75--0.80~\cite{kalpathycramer2016} and confirming the external validity of the perturbation-magnitude calibration described in \S\,2.3.

The 310-nodule cohort was constructed \emph{de novo} under the pre-registered inclusion criteria and is not the conventional 310-nodule benchmark subset used in some prior lung-cancer-classification studies; the 227-patient denominator reflects the present selection cascade and seeded downsampling, not the 110-patient denominator associated with that earlier published subset. All case identifiers in the final cohort carry the \texttt{LIDC-IDRI-} prefix and span the range LIDC-IDRI-0001 to LIDC-IDRI-0601; the full case-ID list is provided as Supplementary File~4.

\subsection{Segmentation perturbation conditions}

For each retained nodule, up to twelve perturbation conditions were generated per paradigm extraction. These comprised: four natural inter-reader masks (R1--R4, with up to three additional masks R5--R7 for the small number of nodules where pylidc clustered more than four annotations), seven controlled synthetic perturbations applied to the per-voxel majority-vote consensus mask $M_C$, and one Dice-calibrated synthetic mask. The synthetic perturbations consisted of morphological erosion with structuring-element radius 1, 2, or 3 voxels (E1, E2, E3); morphological dilation with radius 1 or 2 voxels (D1, D2); and rigid translation by $(\Delta x, \Delta y, \Delta z) \in \{(1,0,0), (0,1,0), (1,1,0)\}$ voxels (T1, T2, T3). A 3-voxel dilation condition was excluded from the protocol after pilot work showed $\sim$40\% volume inflation on small nodules, well outside the empirical inter-reader variability documented in \S\,2.2. All morphological operations were performed in voxel space using SimpleITK 2.3.1 \texttt{BinaryErode} and \texttt{BinaryDilate}; translations used direct array indexing. The Dice-calibrated mask was generated per nodule by adaptive morphological perturbation tuned to match the median pairwise inter-reader Dice for that specific nodule, anchoring the synthetic perturbation magnitude to empirical reader variability rather than to an arbitrary voxel count.

A volume-preservation guard was applied during synthetic-erosion generation: any erosion whose post-operation mask retained less than 25\% of the consensus-mask voxel count was skipped for that nodule, with the omission recorded in the run log. This guard was a necessary consequence of the 6~mm diameter relaxation (\S\,2.2), since fixed-radius erosion can fully obliterate small nodules. Across the 353-nodule extraction cohort, E1 was retained for 48\% of nodules, E2 for 13\%, and E3 for 2\%. The per-nodule perturbation-condition count $K$ therefore varied between 7 and 12, an unbalanced design accommodated explicitly in the intraclass-correlation-coefficient computation (\S\,2.5). Example perturbation masks across the diameter range are shown in Figure~1.

\subsection{Feature extraction}

\subsubsection{Handcrafted IBSI-aligned radiomic features}

Radiomic features were extracted using PyRadiomics 3.1.0~\cite{vangriethuysen2017} in two parallel arms. The three-dimensional arm operated on the full volumetric mask and yielded 107 features per nodule per perturbation, distributed across first-order intensity (18 features), three-dimensional shape (14 features), grey-level co-occurrence matrix (24 features), grey-level run-length matrix (16 features), grey-level size-zone matrix (16 features), neighbourhood grey-tone difference matrix (5 features), and grey-level dependence matrix (14 features) classes. The two-dimensional arm operated on the single maximum-area axial slice of each mask and yielded 102 features per nodule per perturbation, with the 9-feature \texttt{shape2D} class substituted for the 14-feature \texttt{shape} class~\cite{vangriethuysen2017}. The two-dimensional arm was added to enable an apples-to-apples comparison against the foundation-model embeddings, which operate intrinsically on a single slice. Both arms used PyRadiomics default IBSI 1.0~\cite{zwanenburg2020} settings: bin width 25~Hounsfield units, BSpline resampling to $1 \times 1 \times 1$~mm$^3$ isotropic spacing for the 3D arm (the 2D arm preserved native in-plane spacing), and no intensity normalisation. Full parameter files are deposited in Supplementary File~2.

\subsubsection{Foundation-model embeddings}

The foundation-model arm used BiomedCLIP~\cite{zhang2024biomedclip}, specifically the publicly released \texttt{microsoft/BiomedCLIP-Pub\-MedBERT\_256-vit\_base\_patch16\_224} checkpoint. BiomedCLIP is a Vision Transformer image encoder pretrained on 15 million biomedical image-text pairs from PubMed Central via contrastive vision-language learning, with the multimodal projection layer mapping the native 768-dimensional ViT-B/16 output to a 512-dimensional embedding aligned with the PubMedBERT text encoder~\cite{zhang2024biomedclip}. For each perturbation mask, the central axial slice was extracted using the same maximum-area criterion applied in the 2D radiomics arm. The slice was lung-windowed (window level $-600$, window width 1\,500~Hounsfield units), bounding-box-cropped around the mask with 50\% padding, resized to $224 \times 224$ pixels, and passed through the BiomedCLIP encoder. The resulting 512-dimensional embedding was L2-normalised. RadImageNet-ResNet50~\cite{mei2022radimagenet} was planned as a second foundation-model paradigm; weight access was not granted within the study timeframe and this comparison is deferred to a planned amendment. All BiomedCLIP inference ran on CPU, completing in approximately 13~minutes for the 3\,989-embedding extraction.

\subsection{Statistical analysis}

\subsubsection{Intraclass correlation coefficient}

For each feature in each paradigm ($n = 107$ for radiomics-3D, 102 for radiomics-2D, 512 for BiomedCLIP), we computed ICC(2,1) --- two-way random effects, single rater, absolute agreement --- across the universal subset of ten standard perturbation conditions present for every nodule (consensus, R1, R2, R3, D1, D2, T1, T2, T3, DC). The ICC(2,1) variant was selected because both nodule (target) and perturbation condition (rater) are random samples and we wish to generalise beyond the specific readers and perturbations used~\cite{koo2016icc}. Computation used \texttt{pingouin.intraclass\_corr}~\cite{vallat2018pingouin}; ICC(3,1) variants are reported in Supplementary Table~S3 for compatibility with the radiomics QA literature. Features were classified into four standard stability strata~\cite{koo2016icc}: poor (ICC $<$ 0.50), moderate ($0.50 \leq$ ICC $<$ 0.75), good ($0.75 \leq$ ICC $<$ 0.90), and excellent (ICC $\geq$ 0.90).

\subsubsection{Pre-registered hypothesis tests}

\textbf{H1 (proportion-stable comparison):} two-proportion z-test of the fraction of stable features (ICC $\geq 0.75$) between BiomedCLIP and each PyRadiomics arm, Bonferroni-corrected at $\alpha = 0.025$ across the two paradigm-pair comparisons. \textbf{H1' (distribution-shape comparison):} two-sample Kolmogorov--Smirnov test on the empirical ICC distribution per paradigm pair, pre-specified before any data inspection but after observing the stratum-distribution crosstab. \textbf{H2} and \textbf{H3} are downstream-task hypotheses, tested by 5-fold stratified cross-validated L2-regularised logistic regression~\cite{hosmer2013} on the consensus-mask features, with pre- and post-ICC-filter feature sets at thresholds $\tau \in \{0.50, 0.75, 0.85, 0.90\}$ for each paradigm. Mean out-of-fold area under the ROC curve (AUC)~\cite{delong1988auc}, Brier score, and ten-bin expected calibration error were reported with stratified bootstrap 95\% confidence intervals (1\,000 resamples) over pooled out-of-fold predictions. The H3 non-inferiority margin was pre-specified at 0.020 $\Delta$AUC.

\subsubsection{Size-confound mitigation}

To address the structural correlation between nodule diameter and malignancy in LIDC, a diameter-stratified subgroup analysis at strata 6--10~mm, 10--20~mm, and $>20$~mm was pre-specified and conducted with the unfiltered features per paradigm in each stratum, using the same cross-validation and bootstrap protocol described in \S\,2.5.2. The $>20$~mm stratum proved unevaluable owing to LIDC's structural scarcity of large benign nodules ($n \approx 3$ benigns of 76--77 total), a corpus-design limitation discussed in \S\,4.5.

\subsection{Reporting and reproducibility}

The prediction-modelling components of this study were reported in accordance with the TRIPOD+AI 2024 statement~\cite{collins2024tripod}; the foundation-model components followed the Checklist for Artificial Intelligence in Medical Imaging (CLAIM)~\cite{mongan2020claim}; and the per-paradigm risk-of-bias assessment used the Prediction model Risk Of Bias ASsessment Tool (PROBAST)~\cite{wolff2019probast}. All random seeds were fixed at 42 (\texttt{numpy.random.seed}, \texttt{torch.manual\_seed}). Library versions were pinned in a conda \texttt{environment.yml} file; Python 3.9 was required to obtain a Windows-compatible PyRadiomics wheel. Bit-exact reproducibility across two independent runs of the downstream classifier was verified for all 14 active metric rows. The complete code, pinned environment specification, extracted feature tables, per-fold out-of-fold predictions, ICC tables, and figure-generation scripts are released at \url{https://github.com/Iranchamika/paper2-segperturb-fm-radiomics} under MIT licence at the time of manuscript submission.
